\chapter{The register file}\label{ch:register-file}

Each lane of a simdgroup owns a private file of 32-bit general registers, addressed by number in the instruction and
also viewed as 16-bit halves and as 2-, 4- and 8-register tuples. Beside it sits a small uniform file, written once per
dispatch. There is no separate accumulator file for the tensor unit, no register renaming visible to software, and no
interlock that protects a value a program has said it no longer needs. Register sources carry a \emph{modifier} (a few short forms fix it) that can negate it, take its absolute value, or \emph{release} it,
and a release is destructive.

Register allocation on this GPU, and the composition of kernels, must respect three facts:

\begin{enumerate}
\def\labelenumi{\arabic{enumi}.}
\item Fields can name R0–R143, but the hardware executes only R0–R125. An instruction naming a higher register does
  not fault; it is discarded whole, with no error (\cref{sec:namespace}).
\item Fields can be narrower than the namespace: certain forms can name only R0–R15, and others cannot name
  particular registers at all (\cref{sec:narrow-fields-narrow,sec:encoding-holes}). Register choice has to satisfy each form's
  constraints, which a graph colouring over one homogeneous register file does not capture.
\item Lifetime is encoded in the instruction and applies per lane, under the execution mask. The last reader of a value
  marks it released,
  the release destroys the lane's copy (controlled probes then read zero, \cref{sec:release-does}), and no hardware mechanism
  notices a later read (\crefrange{sec:operand-modifiers}{sec:loops-loop-variable-rule}). Liveness is a correctness property of the emitted
  bytes and must be computed per lane over the control-flow graph, including back edges.
\end{enumerate}

When kernels are composed, a value that stays live across a region boundary (a loop header, a fused stage, a
second GEMM) must be kept live in the bytes, inside the 126-register pool, in registers every consuming form can name.

Register notation:

\begin{itemize}
\item \emph{Rn}: general register n, 32 bits per lane. \emph{RnL}, \emph{RnH}: its low and high 16-bit halves. \emph{Ra\_\dots{}\_Rb}: a tuple
  (for example R20\_R21\_R22\_R23, a 16-bit tensor fragment).
\item \emph{MCRegister id}: the number Apple's decoder uses for a register; the general block starts at 105, so Rn has id
  105 + n (verified over 300 registers in 60 shaders; \code{docs/agx3-oracle.md} section 6). The project's code writes
  "r119" for id 119, which is R14; this chapter always writes hardware numbers.
\item \emph{Slot domain}: an operand field value counting halves, 2 × n + h (h = 0 low, 1 high) (decoder, MM~\mm{25.139.1}).
\item \emph{Modifier}: the immediate operand that follows a register source (\cref{sec:operand-modifiers}). \emph{Release} (weight 16), \emph{keep}
  (weight 32).
\item \emph{Late value}: a result that arrives after its producer issues and needs a scoreboard wait (\cref{sec:scoreboard}).
\item \emph{Registers used}: the highest 32-bit register named in the main program, plus one; it equals the declared count
  this compiler writes into metadata slot 0 (\cref{sec:gpu-metadata-slots}).
\item The running example (\cref{sec:one-program-ir}) is used in \cref{sec:release-does}.
\end{itemize}

\Cref{tab:register-statements} lists the chapter's main statements and their standing.

\begin{xltabular}{\linewidth}{@{}XX@{}}
\caption{Main statements of the chapter and their standing.}\label{tab:register-statements}\\
\toprule
Statement & Standing \\
\midrule
\endfirsthead
\tablecontinued{2}
\toprule
Statement & Standing \\
\midrule
\endhead
R0–R125 execute; naming R126–R143 squashes the whole instruction & measured (MM~\mm{10}, MM~\mm{25.18}) \\*
Release destroys the lane's value, per lane, gated by
the execution mask & measured (MM~\mm{6}, MM~\mm{25.95}) \\*
A released register's later read returns 0 & observed in every controlled run; not established as a rule (
\cref{sec:release-does}) \\*
Loop-carried values must not be released before the back edge & measured, with a machine-hanging failure (MM~\mm{25.116}) \\*
8N + K ≤ 115 and 8N + 4M ≤ 108 & measured frontiers of the kernels tested (MM~\mm{11}) \\*
Field widths and holes & decoder (Apple's decoder reading back the
encoder's bytes) \\*
\bottomrule
\end{xltabular}

\section{Register namespace}\label{sec:namespace}

Allocate only R0–R125: 126 registers per lane, fifteen eight-register groups at bases R0, R8, \dots{}, R112
(measured, MM~\mm{10}, MM~\mm{0.7}). Any instruction that names R126, R127 or R128–R143 in \emph{any} register operand is squashed as a
whole: no register write, no memory write, no atomic, no tensor-unit time, and no error from the compiler, the driver
or the API (measured, MM~\mm{17} rule 1). Nothing faults, and a stray high operand shows up only as wrong output, so the
allocator has to enforce the bound when it assigns registers.

The boundary first showed up as a store of a tuple reaching R126 that failed to round-trip, independent of the
declared count. A follow-up wrote through R126 and R127 and read zeros from the output buffer, and was recorded as
R126 and R127 \emph{reading as zero and dropping writes}, like a hardwired zero register. The rival reading is that the
hardware discards any instruction that names them. The two predict different results for an instruction that reads
R126 and stores it: the first stores 0, the second stores nothing. That experiment could not separate them, because the
output buffer had been zero-filled before the dispatch and "stored 0" and "stored nothing" leave the same bytes. With
the buffer filled with a sentinel instead, the sentinel survived, so the store had not executed (MM~\mm{19}, MM~\mm{25.20}).
The same test was then repeated in every operand role: ALU sources and destinations, store sources of three widths, a
store's index register, a load's destination group, an atomic's destination, and the tensor MMA's accumulator, A and B
operands (MM~\mm{10}). In each the whole instruction is discarded, not only the register access. For R128–R143 the result
is unchanged by declared register counts from 12 to 255, by one against 32 simdgroups, and by a 32- against a
1,024-thread pipeline limit (MM~\mm{25.18}). The wide RMSNorm's first check (\cref{sec:release-does}), which passed against an output that
carried no sentinel, had the same flaw.

\code{agxforge/g17/registerdomain.py} encodes the rule: \code{ALLOCATABLE\_MAX = 125}, \code{REGISTER\_COUNT = 126},
\code{MAX\_FP32\_ACCUMULATOR\_GROUPS = 15}, group starts on a multiple of 8 (this compiler). The boundary is measured on
one part; whether another G17 implementation backs the higher names with storage is not tested (MM~\mm{25.18}).

Beyond R125, the encoding and the decoder distinguish four cases.

\begin{itemize}
\item \emph{R126 and R127} are the top pair of a 128-register architectural file that this part populates only to R125. Three
  routes agree on 128: the half-register move's seven-bit source field plus a one-bit file selector (128 halves × 2
  files = 128 registers, measured in this repository); the five-bit accumulator selector's eighteen groups, of which
  exactly fifteen lie inside R0–R125, one straddles (R120–R127) and two lie wholly beyond; and, as a hypothesis source
  only, Mesa's G13 model of 256 sixteen-bit slots. Sentinel stores rule out a hardwired zero and unique encodings rule
  out aliases, and no metadata or launch control enables the two registers (MM~\mm{24.2} row H1, MM~\mm{25.20}). Whether anything \emph{reserves} the
  top pair, rather than the part simply populating 126 of 128, is open; only documentation could settle it.
\item \emph{R128–R143} are decoder-visible names outside that file, reachable only through fields wide enough to express them
  (accumulator-selector groups 16 and 17, nine-bit slot fields). Load and store register fields reach R127 only, so for
  those forms R126/R127 are expressible and R128–R143 are not (MM~\mm{24.2} row H1, MM~\mm{25.38}). Prior art's larger uniform file
  is selected by operand type, not by a high register number, so it does not explain them (MM~\mm{25.18}).
\item \emph{The encoding} admits R0–R143 and rejects 144–255 (MM~\mm{24.2} row H2).
\item \emph{Accumulator selector values 18–31} decode to no operand at all: null targets with no name, absent from every
  compiled witness, never dispatched because malformed tensor encodings have hung the GPU (decoder, MM~\mm{25.19}). An
  assembler must reject them.
\end{itemize}

\Cref{tab:highest-base} gives the highest legal base register for each value class (measured, MM~\mm{0.7}, MM~\mm{23.2}).

\begin{xltabular}{\linewidth}{@{}Xlll@{}}
\caption{Highest legal base register per value class.}\label{tab:highest-base}\\
\toprule
Value & Registers & Alignment & Highest base \\
\midrule
\endfirsthead
\tablecontinued{4}
\toprule
Value & Registers & Alignment & Highest base \\
\midrule
\endhead
scalar 32-bit & 1 & 1 & R125 \\*
int8 / uint8 tensor fragment & 2 & 2 & R124 \\*
half / bfloat tensor fragment & 4 & 4 & R120 \\*
fp32 tensor fragment, fp32 / int32 accumulator & 8 & 8 & R112 \\*
\bottomrule
\end{xltabular}

\section{Register classes}\label{sec:register-classes}

Apple's decoder tables survive with register and register-class names (3,567 registers, 500 classes); instruction
names do not (decoder, \code{docs/agx3-oracle.md} section 6). \Cref{tab:register-classes} lists the classes used in this document, with member counts read from those
tables (\code{tools/agx3meta classes}, \code{members}).

\begin{xltabular}{\linewidth}{@{}>{\hsize=0.72\hsize}Xll>{\hsize=1.28\hsize}X@{}}
\caption{Register classes in Apple's decoder tables, with width and member count.}\label{tab:register-classes}\\
\toprule
Class & Width & Members & What it is \\
\midrule
\endfirsthead
\tablecontinued{4}
\toprule
Class & Width & Members & What it is \\
\midrule
\endhead
GPR32 & 32 & 144 & R0–R143 \\*
GPR16 & 16 & 288 & R0L, R0H, \dots{}, R143H \\*
GPR32tup2 / tup4 & 64 / 128 & 143 / 141 & 2- and 4-register tuples \\*
GPR32tup8\_alignedrc & 256 & 18 & the eighteen 8-aligned groups R0–R7 \dots{} R136–R143 \\*
IR32 & 32 & 16 & IR0–IR15, MCRegister ids 89–104 \\*
IRGPR32 & 32 & 160 & R0–R143 followed by IR0–IR15 \\*
FLAGR / FLAGwritable & 16 & 17 / 15 & FLAG0–FLAG14 plus the constant sources FLAGTRUE, FLAGFALSE (
\cref{ch:control-flow-execution}) \\*
\bottomrule
\end{xltabular}

A 16-bit operand names a half of a general register; the slot domain numbers halves 2n + h, confirmed
through Apple's decoder (MM~\mm{25.139.1}). Writes and releases act per half (\cref{sec:release-does}): a 16-bit write leaves the other
half's contents and history untouched. Measured for \op{14060}: with R holding
7t + 0xABCD0000, writing its low half with the lane index left 0xABCD0000 \textbar{} lane, three runs identical; it does not
zero-extend (MM~\mm{25.145.4}).

Register fields are encoded in units of the tuple width, so a tuple base must be naturally
aligned (2, 4 or 8); an unaligned request cannot be expressed and must be refused, not rounded (MM~\mm{0.4}, MM~\mm{2}; the
encoder bug that rounded silently is \cref{sec:admission-decodes-executes}). The tensor unit's operands use three classes found nowhere else in
the table: class 335 (8-register tuple, the destination of every 32-bit-accumulator form), class 159 (4-register
tuple) and class 41 (2-register tuple); there is no fourth, hence no sub-byte fragment class (decoder, MM~\mm{25.6}). Eight-bit
slot fields come in two addressing conventions (half-indexed: the value 128 names R64, on 3,970 operands;
register-indexed: 128 names R128, on 638), and one scalar form carries one of each, so field width alone does not say
which registers a field reaches (decoder, MM~\mm{25.38}).

The IR file proper is the 16-member class IR32. IRGPR32 is a class that \emph{admits} the IR
registers alongside the general ones (R0–R143 followed by IR0–IR15, 160 members); its operands decode as ordinary
registers, value 105 = R0, and all 229 corpus writes by its users (\op{10239}, \op{10824}, \op{11060}, \op{14147}) land in Rn that ordinary instructions read (MM~\mm{25.59}). The corpus uses only its general members. What the 16
IR registers hold, and whether any executed program names one, is not recorded (open).

The uniform file is a per-dispatch file of 32-bit uniforms u0, u1, \dots{}, written by the \emph{constant program} (a second
program Apple's compiler emits, run once per dispatch before main; \cref{sec:constant-program-record}) and read by main-program sources.
Apple's decoder writes a uniform operand as \code{bin(op0, const(K), S)}, where op0 is the \op{0}, and K counts 16-bit halves, so the uniform is K // 2; a 64-bit value occupies u(2k), u(2k+1) (decoder, MM~\mm{25.120}). In tensor (MMA)
programs u0–u5 are reserved and constant-program writes start at u6; outside MMA programs some constant programs write
below u6 (Apple's compiler emits, MM~\mm{25.120}). A uniform destination must be written with one value across the active
lanes (measured by emulation against Apple's GPU, MM~\mm{25.179}). Uniforms also serve as operands of ordinary forms: the mask
operand of \op{428} in its pool form is a uniform halfword preloaded from the slot-13
pool (MM~\mm{25.141.17}, MM~\mm{25.145}; \cref{sec:constants-uniforms}).

\Op{592} stages operands for texture messages in the \op{4}; the writer names a scoreboard slot and the message instruction waits on it
(decoder over 189,898 staging writes, MM~\mm{25.120.1}; \cref{sec:textures}, \cref{sec:scoreboard-publish-wait}).

FLAG0–FLAG14 are an allocatable predicate resource with pressure, not one condition code; compact forms
reach FLAG0–FLAG6 only (decoder, \code{isa/g17-special-registers.toml}). They are covered in \cref{ch:control-flow-execution}.

\section{Narrow register fields}\label{sec:narrow-fields-narrow}

Each operand field has its own width, and the allocator must compute, per value, the highest register
every form touching it can name. The authoring table records field widths per (opcode, length) form, and this
compiler derives each value's ceiling from them (\code{cc.\_caps}), with a hand list only for forms outside the table
(this compiler, \code{agxforge/g17/cc.py} \code{\_needs\_narrow}). The binding narrow field today is the four-byte special-register
read's four-bit destination, R0–R15 (decoder and this compiler, MM~\mm{25.110.1}, \code{agxforge/g17/cc.py}). Other forms carry
field-specific ceilings: the fp16 conversions author destinations to R61 with holes (\cref{sec:encoding-holes}), and in measured
kernels the fma's first source and a load's index register were confined to narrow registers by their forms' fields at
the time (MM~\mm{25.144.3}).

This compiler keeps R4–R15 (twelve registers) as its narrow pool and places a value that is only
ever an ALU first source in the wide pool (R16 and up) first, to keep the narrow registers free (this compiler,
MM~\mm{25.146}). For one witness the policy makes this compiler declare 17 registers where Apple's compiler
declares 3; which policy gives better occupancy is unmeasured (MM~\mm{25.146}). Refusals name the pool ("12 narrow + 110
wide available"); when they occur, the narrow pool, not the file, ran out (MM~\mm{25.144.3}).

The widths now used were measured after the early scalar field map proved wrong. That map read the ALU
destination as five bits and the load, store, special-register-read and move-immediate destinations as four bits in
byte 0, which implied a 32-register file (the ALU destination saturated at 31 in kernels needing 32 to 64 live values).
Apple's own shaders contradicted that (39 percent exceed R31; driver shaders reach R112). The ceiling was then measured on
the hardware, with one value pinned into one high register and read back, one dispatch per process, with the expected value 107
preregistered every time. R20, R40, R64, R100, R120, R124 and R125 read back correctly, and R126 and above were squashed (measured, \code{agxforge/g17/cc.py} comment,
\code{spike/accel/re/hireg.py}). The fields were re-measured at seven and eight bits (extension bits in byte 2 and byte 5).
The hand list also hid one under-measured field: a register-register bitwise form that it held to the narrow
registers had a destination field recorded as five bits. The field was six bits, which went unnoticed until something
was allocated above R15 (this compiler, \code{agxforge/g17/cc.py}).

\section{Operand modifiers}\label{sec:operand-modifiers}

Every register source is followed by a modifier immediate. Over 12,000 corpus instances of five opcodes
it takes exactly the values \{0, 2, 16, 18, 20, 32, 34, 36\}: bit 1 (2) negate, bit 2 (4) absolute value, bit 4 (16)
release, bit 5 (32) keep, composing freely (18 = negate + release, 34 = negate + keep) (Apple's compiler emits;
\code{agxforge/g17/auth.py}, \code{isa/g17-source-modifier-operand.toml}). A destination's modifier takes only 0 or 32: a definition
is always a keep. Value 0 carries no lifetime information and preserves the register (the running example's first
reads, \cref{sec:release-does}, execute exactly with 0). Floating subtraction is \op{998} with a negate modifier on the
second source; \code{a - b} and \code{a + (-b)} compile to the same bytes (Apple's compiler emits, MM~\mm{25.146}).

In the ALU and tensor forms the destination's raw operand (operand 1, which
the decoder prints as one wide integer and older notes call the "flags word") packs several fields: bits 4–5 the
destination's lifetime, bits 20–23 a fill slot, bits 24–31 the scoreboard wait mask (measured, MM~\mm{25.145}, MM~\mm{6}; the fields are \cref{sec:scoreboard-publish-wait}). In the
tensor MMA the keep bit is operand 1 weight 32, physically byte 4 bit 3 (\code{auth.carriers(5106, 1)[5] = [(4, 3, 0)]}), and
it is hardware-inert: a conservative assertion that the next MMA writes the same destination tuple (\cref{sec:modifiers-keep-flag}).

MM~\mm{6} and \code{agxforge/g17/auth.py} give bit 5 as keep; MM~\mm{25.59} reports that
\op{10239}'s 101 instances carry "source modifier 16 (32 for a high-half source)" in unsigned
kernels and 24 (40) in signed ones. These are consistent: in
the saturating add the modifier carries an extra bit 3 (weight 8) that selects signed saturation (signed arm executed and clamped to
INT\_MIN..INT\_MAX, MM~\mm{25.59}), so the field's low bits are not uniform across forms. The 16 / 32 split itself is
consistent with release and keep: 16 + 8 = 24 and 32 + 8 = 40 preserve it. MM~\mm{25.59} records a correlation over 101
instances, not a measurement isolating the bit, and nothing shows a high-half \emph{selector} at weight 32. A compiler may
treat 32 as keep in the saturating add as in the general census; it may not assume that every form's modifier is limited to the
four general bits.

In the forms of \cref{tab:fixed-modifiers}, the modifier cannot be chosen freely.

\begin{xltabular}{\linewidth}{@{}>{\hsize=0.95\hsize}X>{\hsize=1.05\hsize}X>{\hsize=1.05\hsize}X>{\hsize=0.95\hsize}X@{}}
\caption{Forms whose modifier cannot be chosen freely, with the consequence for a compiler. Each row is measured or decoder-checked.}\label{tab:fixed-modifiers}\\
\toprule
Form & Fact & Consequence & Source \\
\midrule
\endfirsthead
\tablecontinued{4}
\toprule
Form & Fact & Consequence & Source \\
\midrule
\endhead
four-byte \op{424}, \op{13575}, \op{17771} & no lifetime bit: always release their sources & use the longer form for a source read again & measured, MM~\mm{24.8} row T4, ledger
\code{g17-the-four-byte-bitwise-releases-its-sources} \\*
\op{10094}, ten bytes & the value's lifetime (operand 9) is pinned at 16 & two atomics sharing a value, or the value stored afterwards, read 0
(fuzz seed 5012: 0 where 24) & measured, MM~\mm{25.144.6} \\*
\op{17256} & the stored tuple's lifetime (operand 1) is 16 in every
witness; no keep value measured & a value read after the store read 0 on 32 of 32 lanes; store a copy & measured, MM~\mm{25.144.6} \\*
\op{14157} & the source lifetime is operand 3, not operand 1; a
template's 16 released x & \code{y = broadcast(x); z = fma(x, x, y)} stored 0 & measured, MM~\mm{25.144.6} \\*
imageblock store & operand 6 is the coordinate's lifetime (0 keep, 16
release) & releasing it collapsed all 32 lanes to element 0 (one-bit control) & measured, MM~\mm{6} \\*
\op{5106} & the accumulator's modifier (operand 10) has no
single-bit carrier in the 80-bit field map & cannot be authored; a multi-bit encoding is not excluded & decoder, MM~\mm{6} \\*
\bottomrule
\end{xltabular}

The position of the lifetime operand varies by form (\op{586}: operand 3; \op{17229}/8: operands 1 and 6 for value and index; the fp16 fma, register form: operands 3, 5, 7), and each must be
written from liveness (MM~\mm{25.117}, MM~\mm{25.196}, \code{agxforge/g17/cc.py}). All four modifier bits must be written, not only the
lifetime: witnesses taken from a mutation walk rather than from Apple's code can carry 0 or 52 (keep, release and
absolute value at once), and inheriting them made four opcodes in a naming sweep return the same value for every input
(\code{agxforge/g17/auth.py}).

\section{Release semantics}\label{sec:release-does}

A release (modifier 16) destroys the lane's value, per lane, under the execution mask, and a
later write restores the register. Every controlled probe of a released register read zero, but that is not
established as a rule (see the end of this section).

Each of the following was measured in executed programs.

\begin{itemize}
\item The contents do not survive, in the controlled probes. An instruction carrying 16 as the first reader of a
  register made every later reader in that program see 0, and the register read 0 at a store. Modifier 32, or another
  first-reader form, preserved it (MM~\mm{6}). The imageblock store's coordinate behaved the same way, so the effect is not
  specific to the tensor feed (MM~\mm{6}).
\item The release acts per lane, gated by the execution mask. Every earlier release had run on all 32 lanes, so two
  mechanisms fit equally well: the release frees a register the simdgroup shares, or it clears each lane's copy. A
  partial mask separates them. \code{y = t + 3054} was read by \op{11666}/12 inside a
  region gated on \code{t < bound}, with that read's lifetime written 32 → 16 and decoded back; after the region every lane
  stored \code{y + 0} (MM~\mm{25.95}, \code{tools/g17releasemask.py}). \Cref{tab:release-mask-arms} gives the four arms.

  \begin{xltabular}{\linewidth}{@{}lllX@{}}
\caption{Lanes reading 0 after a release under each execution mask (MM~\mm{25.95}).}\label{tab:release-mask-arms}\\
\toprule
Arm & Region & Patched & Lanes reading 0 \\
\midrule
\endfirsthead
\tablecontinued{4}
\toprule
Arm & Region & Patched & Lanes reading 0 \\
\midrule
\endhead
  control & t < 16 & no & none \\*
  partial & t < 16 & yes & 0–15 only; 16–31 read y exactly \\*
  all on & t < 32 & yes & all 32 \\*
  all off & t < 0 & yes & none \\*
  \bottomrule
\end{xltabular}

  In the partial arm only the lanes that executed the release lost y, and in the all-off arm none did: the release acts
  on each lane's copy, an instruction no lane executes releases nothing, and a simdgroup-wide free is ruled out. A
  compiler must therefore treat a value released inside a branch as dead only on the lanes that took the branch
  (\cref{sec:loops-loop-variable-rule}, rule 3).
\item The register is cleared, not handed to another simdgroup, within the measured bound. Group 0 released y (one
  lifetime operand, the only bytes that differ from the keep arm), stayed resident for about 13,100
  instructions while 8,191 other groups (about 300 instructions each, 16 distinctive live values each) launched and
  finished, then re-read y: 0 in 3 of 3 rounds, never another simdgroup's value, while group 1 (not released) read y in
  the same dispatch (MM~\mm{25.95}, \code{isa/g17-execution-release-reuse-results.json}). This bounds when reuse could occur but
  does not exclude reuse.
\item A write restores the register: 230 release → rewrite → read chains in 15 bit-exact tensor bodies, rewritten by ALU results,
  immediate moves, loads and a tensor MMA; for example \op{998} releases R19, \op{3290} writes R19,
  and the add reads it (MM~\mm{25.95}).
\item Order inside one instruction depends on the form. fp32 multiply squaring \code{x * x} with a release on the first read
  is bit-exact, and \op{9700} naming one register in two released slots is harmless: in
  these ALU forms the release takes effect after the sources are read (MM~\mm{25.145}, MM~\mm{25.117}). But \op{17229} naming one register as both value and index, with the \emph{index} slot releasing it, stores 0 on every lane: the
  store reads its value after the index's release has cleared it (\code{5f04030a01061040} decodes to {[}R5, 16, 2066, expr,
  0, R5, 16{]}; clearing operand 6 alone restores the value). Apple's compiler releases neither slot in all 142 such
  stores among 32,162 (MM~\mm{25.117}). Rule: never release, in one slot, a register another slot of the same instruction
  reads, except in forms measured safe.
\item A write cancels the destination's pending release. An in-place accumulate that reads C with a release and writes
  D = C + \dots{} with D = C is bit-exact (MM~\mm{25.145}).
\item Release acts per 16-bit half. A 16-bit write revives only its half; the other keeps its history (MM~\mm{25.145}).
\end{itemize}

Whether a released register reads 0, or can still return its old value, is unresolved (the sources' differing wordings
are listed in \cref{app:a}). The evidence on each side is as follows.

\begin{itemize}
\item every controlled measurement read 0: adjacent reads, one ALU op later, about 7 instructions later, past a barrier,
  under partial masks, after 13,100 instructions of residency (MM~\mm{25.95}); and the M 2048 register-O attention whose
  32-bit read of R33 needed its released high half to be zero, which passed 20 of 20 in fresh processes once
  the 16-bit special-register read was shown not to zero-extend (\cref{sec:register-classes}; MM~\mm{25.147}, MM~\mm{25.145.4});
\item the only "old value" observation is a loop kernel (the wide RMSNorm's scale loop, \cref{sec:loops-loop-variable-rule}) that passed one check
  and then failed with the same bytes; the check that passed had no sentinel in the output, so leftover correct values
  could have hidden unwritten ones (MM~\mm{25.142.7}, \code{agxforge/g17/releasecheck.py} header).
\end{itemize}

The old-value branch is weakly evidenced but not excluded. A compiler may assume that a released register never
shows another simdgroup's value within the measured bound, and that rewriting it restores it. It may not assume it
reads 0, and may not assume it keeps its value. It may read one only where both outcomes give the same answer, which
is the rule this project's emulator applies (it admits such a read only when the old value is zero, MM~\mm{25.145}) and the
rule the attention loop's latch relies on (the counter a finished lane released is the zero that ended its loop).
\emph{Correction:} MM~\mm{0.15}'s "A released register: reads zero" and MM~\mm{17} rule 37's "every later read returns zero" state an
observation as a rule.

Of the never-written registers, untouched R2, R20 and R60 read zero as operands in a whole-program measurement; R0 does not
(measured, MM~\mm{25.145}; the emulator admits such reads only in its whole-program tier). Do not rely on R0.

Apple's compiler reads released registers in one place. Its legacy simdgroup MMA programs release R0 and R1
(address arithmetic) and then issue \op{2842} with accumulator R0\_R1 (Apple's compiler
emits, \code{agxforge/g17/releasecheck.py}). The checker treats this as a zero initialisation and exempts it; it is a statement
about what Apple emits, not a measurement of the read. The checker treats Apple's 6,594 corpus programs as ground
truth and is required to report zero findings on them; this exemption, one for cross-lane value operands and a few
path-sensitivity rules were each added for a named false positive there (\code{agxforge/g17/releasecheck.py}).

\textbf{Worked example: release at the last read in the running example's four MMAs} (compiled by this project's compiler;
decoded by Apple's decoder; executed below Metal, 323 of 323 exact). The product is 17 × 19 × 16, so the M and N axes
each need two 16-wide tiles (the second masked) and K = 16 needs one step: four \op{5107} instructions, each reading one A tile and one B tile. Each tile is read twice, so its first read
carries modifier 0 and its second, last read carries 16 (\cref{tab:example-mma-releases}).

\begin{xltabular}{\linewidth}{@{}lX>{\hsize=1.21\hsize}X>{\hsize=0.91\hsize}X>{\hsize=0.91\hsize}X@{}}
\caption{Tile modifiers in the running example's four tensor MMAs, with the release on each tile's last read.}\label{tab:example-mma-releases}\\
\toprule
Offset & Bytes & D (accumulator tile) & A tile, modifier & B tile, modifier \\
\midrule
\endfirsthead
\tablecontinued{5}
\toprule
Offset & Bytes & D (accumulator tile) & A tile, modifier & B tile, modifier \\
\midrule
\endhead
0x29a & \hash{a70405112204a4620202} & R24–R31 (rows 0–15, cols 0–15) & R20\_R23, \textbf{0} & R68\_R71, \textbf{0} \\*
0x2a4 & \hash{2704a5112204a4420204} & R32–R39 (rows 0–15, cols 16–31) & R20\_R23, \textbf{16} (last read) & R72\_R75, \textbf{0} \\*
0x2ae & \hash{a70485832206a4420202} & R40–R47 (rows 16–31, cols 0–15) & R56\_R59, \textbf{0} & R68\_R71, \textbf{16} (last read) \\*
0x2b8 & \hash{2704a5832206a4620204} & R48–R55 (rows 16–31, cols 16–31) & R56\_R59, \textbf{16} (last read) & R72\_R75, \textbf{16} (last read) \\*
\bottomrule
\end{xltabular}

R20\_R23 holds A rows 0–15 (two tensor loads at byte offsets 0 and 256, eight rows of 32 bytes each); R56\_R59 holds A
rows 16–31, of which only row 16 exists (masked loads at 512 and 768); R68\_R71 and R72\_R75 hold B columns 0–15 and 16–31
(row stride 38 bytes, so offsets 0 / 304 and 32 / 336). The A modifier is operand 4, carried by byte1 bit 7 (weight 32)
and byte2 bit 5 (weight 16); the B modifier is operand 7, carried by byte5 bit 1 (16) and byte8 bit 0 (32)
(\code{isa/g17-tensor-mma-fieldmap.json}). The first two instructions differ in four bytes: byte 2 bit 5 sets A's release
(0 → 16); D moves from R24 to R32 through byte 0 bit 7, byte 2 bit 7 and byte 7 bit 5 (weights 8, 32 and 16 on the
tuple's base register); byte 9 moves B from R68 to R72. Nothing reads R20–R23, R56–R59 or R68–R75 afterwards; each D
tile is read once, as two four-register halves, by the tensor stores that follow.

Placing the release on the \emph{first} read of any tile would have zeroed one of the four products on every lane. Three properties
of the schedule make the placement trivial here and are what a general allocator must establish: the four MMAs run on
every lane (no mask divides them), there is no loop (no back edge reads a tile again), and no instruction reads a tile in
two slots.

\section{Loop-carried values}\label{sec:loops-loop-variable-rule}

A loop compiled by this compiler (measured, MM~\mm{25.116}), bounded by a register
(\code{i < trips}) and capped by a constant (\code{i < 255}) so that it was "bounded on paper", never returned. The host was killed
at its 600 s timeout and the kernel kept running: with nothing dispatching, the GPU read 100 percent device utilisation
until a reboot. Decoded offline, the latch was \code{icmp(i1 < trips)} followed by the cap
\code{icmp(i1 < 255)}, and the cap compare, the instruction added to make termination provable, carried a release on the
counter. The next trip read the cleared counter, \code{i} stuck at 1, and neither bound could ever be reached.
The compiler's liveness took a value's last textual read as its last live point; a loop header's phi carries no uses, so
the latch value, read only inside the body, looked dead after its last mention. Every earlier runtime-loop test stored its
counter after the loop, which kept it live and hid the defect.

A loop phi that inherits its entry value's register must keep
that value live into the loop (measured, MM~\mm{25.142.7}, MM~\mm{25.144.6}). In the wide RMSNorm, the scale loop's index phi started at the thread index t and took
t's register; t's last mention before the loop released it, so every thread stored as thread 0 (only out{[}0{]} and out{[}1024{]}
written). This holds even when the instruction falling into the header itself reads the entry value: the value's last
point is the loop header label, because that instruction would otherwise be its last reader by index (\code{c = B[t] + 7; x = c + 1; loop i = phi(c, i + 1)} returned 2 instead of 1009 on 32 of 32 lanes; 29 of 150 fuzzed loop programs had the pattern, all with
phis starting at constant 0, which a released register also read, so they were right by luck).

This compiler applies three liveness rules.

\begin{enumerate}
\def\labelenumi{\arabic{enumi}.}
\item A back edge \emph{reads} its header phis' latch values; liveness must treat it so (this compiler: \code{Alloc.\_cfg\_uses}).
\item Refuse a release of any loop-carried register (one read in the body before the body writes it) after the body's last
  write of it (MM~\mm{25.116}).
\item Compute releases per lane over every path a lane can take, with masks and loops, to a fixed point: a value released
  inside a branch is dead only on the lanes that took it, so it may not be treated as dead after the join (MM~\mm{25.95}).
  Liveness taken from the last mention in instruction order is wrong for loops, where a back edge reads a value after its last mention.
\end{enumerate}

\code{agxforge/g17/releasecheck.py} implements rule 3 on emitted bytes, walking each lane's paths with the execution-mask
counter and joining released-half sets by union to a fixed point; it runs after every compile, held to zero findings on
Apple's corpus (this compiler, MM~\mm{25.144.6}; \cref{sec:release-does}). The check has to be static, because the looping
kernel at the start of this section could not be stopped from the host (MM~\mm{25.116}). Because a released read can pass a
check once, loop kernels must also be verified against sentinel-filled outputs (MM~\mm{17} rule 37).

\section{Encoding holes}\label{sec:encoding-holes}

Some forms cannot name particular registers: no byte pattern the encoder knows decodes back to them. These are found
by authoring each register and reading it back through Apple's decoder (decoder only); \cref{tab:encoding-holes} lists them.

\begin{xltabular}{\linewidth}{@{}>{\hsize=1.00\hsize}X>{\hsize=0.85\hsize}X>{\hsize=1.40\hsize}X>{\hsize=0.75\hsize}X@{}}
\caption{Registers each holed form cannot author, found by reading each authored register back through Apple's decoder.}\label{tab:encoding-holes}\\
\toprule
Form & Operand & Registers it cannot author & Source \\
\midrule
\endfirsthead
\tablecontinued{4}
\toprule
Form & Operand & Registers it cannot author & Source \\
\midrule
\endhead
\op{1004}/12 & destination & R14, R15, R30, R31, R46, R47, and R62 and above (the encoder reads back
d + 2, or refuses a forced bit) & MM~\mm{25.133}, \code{agxforge/g17/cc.py} \\*
\op{798}/12 & low-half destination & R18, R19, R26, R27, R50, R51, R58, R59 (low halves) & MM~\mm{25.196} \\*
\op{798}/12 & operand 4 / operand 6 sources & R17L / R18L & MM~\mm{25.196} \\*
\bottomrule
\end{xltabular}

\emph{Correction:} the fp16 fma's holes are the listed set, not "registers 2 and 3 mod 8"; MM~\mm{25.136.6}'s widening holes
(d = 14 and 15) are a subset of the set above (\cref{app:a}).

To avoid the holes, this compiler's allocator sweeps the admissible indices once from the encoder rather than typing them in
(\code{\_encodable\_dests}, \code{\_encodable\_fma16\_sources}), and offers them only to values those forms define or read. A third,
hole-aware allocation attempt, used only after two refusals so earlier programs keep their bytes, gives unconstrained
values the holed forms' unauthorable registers first (MM~\mm{25.136.6}). The holes surfaced as compile crashes in a fuzzer
(seed 6) and as refusals of loop arms, never as wrong bytes, because every emitted operand is read back (\cref{sec:admission-decodes-executes}).

Before per-length slot read-back existed, two forms misencoded silently, without raising:
\op{11994}/16's operand 5 named R105 for every register from R16 up, and \op{3291} and two sibling forms at four bytes wrapped R64 and above by writing past byte 4 (decoder, MM~\mm{25.88}). The
encoder now reads every slot back through the decoder-measured layout for its length and refuses a value it cannot hold.

\textbf{Not known.} Whether these registers are unreachable in the encoding or only in the project's field maps (a carrier
bit Apple's instances never vary would be absent from the maps).

\section{Register budgets and spills}\label{sec:capacity-budgets-spills}

The capacity rules are measured (MM~\mm{11}, MM~\mm{17} rules 2 and 3).

\begin{itemize}
\item One flat pool of 126 registers: accumulators, fragments and scalars trade capacity; there is no separate accumulator
  resource and no aliasing. Widths: fp32 accumulator 8, fp32 operand fragment 8, half or bfloat fragment 4, int8 operand
  pair 2, scalar 1.
\item At most 15 live eight-register fp32 accumulator groups in straight-line code (122 registers touched; 16 spill 24
  stores), 14 in a loop with one A and one B fragment live.
\item With one A and one B fragment live, the measured trade lines are \code{8N + K <= 115} for N accumulators against K live
  scalars and \code{8N + 4M <= 108} against M 16-bit fragments, exact at every point tested. The constants include the
  measured kernels' own overhead registers; they are frontiers of those kernels, and the hardware bound is the 126-register
  pool.
\item Spill points in Apple's compiler are scheduling outcomes, where its backend schedule overflows the pool, not
  register-file boundaries (peaks at the cliffs 118–125). Count both spill forms when diagnosing: \op{595} and \op{17789}; a scalar-heavy kernel had 0 of the first and
  74 of the second (Apple's compiler emits, MM~\mm{0.7}, MM~\mm{11}, MM~\mm{13}).
\item This compiler does not spill: it refuses a program that does not fit, naming the pool that ran out (this compiler,
  MM~\mm{25.113}, MM~\mm{25.144.3}).
\end{itemize}

What spills, registers and occupancy cost in time is measured in \cref{sec:registers-spills-occupancy,sec:residency-occupancy-measurements}; residency against registers
used is a rule in \cref{sec:residency-occupancy-rules}.

\textbf{Evidence.} MM~\mm{0.4}, MM~\mm{0.7}, MM~\mm{0.15}, MM~\mm{2}, MM~\mm{6}, MM~\mm{10}, MM~\mm{11}, MM~\mm{13}, MM~\mm{17} (rules 1–3, 37), \mm{19}, \mm{23.2},
\mm{24.2} (rows H1, H2), \mm{24.8}, \mm{25.6}, \mm{25.18}, \mm{25.19}, \mm{25.20}, \mm{25.38}, \mm{25.59}, \mm{25.88}, \mm{25.95},
\mm{25.110.1}, \mm{25.113}, \mm{25.116}, \mm{25.117}, \mm{25.120}, \mm{25.133}, \mm{25.136.6}, \mm{25.139.1}, \mm{25.141.17},
\mm{25.142.7}, \mm{25.144.3}, \mm{25.144.6}, \mm{25.145}, \mm{25.145.4}, \mm{25.146}, \mm{25.147}, \mm{25.179}, \mm{25.196}; \code{agxforge/g17/releasecheck.py},
\code{agxforge/g17/cc.py}, \code{agxforge/g17/auth.py}, \code{agxforge/g17/registerdomain.py}, \code{isa/g17-tensor-mma-fieldmap.json}.
